\section{Lean Definitions and Theorems}
\label{app:lean}

This appendix lists the Lean~4 statements behind the results marked
\textsuperscript{\textsf{\tiny Lean}}, grouped by section, with proofs
omitted. Names are those of the development. The type
\lstinline{ℝ≥0∞} is $\ENNReal$, a formula of
Section~\ref{sec:linear} is a term of the inductive type \lstinline{Fm},
\lstinline{web A} is its web, \lstinline{P A} its probabilistic coherence
space, \lstinline{recipe A} its calculus, and \lstinline{loc A} its local
calculus. Every listed theorem depends only on the axioms
\lstinline{propext}, \lstinline{Classical.choice}, and
\lstinline{Quot.sound}, except the choice-free results, which avoid
\lstinline{Classical.choice}.

\subsection{The Development}

The development builds with \textsf{mathlib} and contains no
\texttt{sorry}. Every listed theorem depends only on the axioms above,
which \texttt{\#print axioms} confirms. The files that contain the
results of this paper are the following, with their sizes in lines.

\begin{center}
\small
\begin{tabular}{@{}lrl@{}}
\toprule
File & Lines & Content \\
\midrule
\texttt{Recipe} & 603 & presentations, mixtures, affine laws, completeness \\
\texttt{Farkas} & 648 & Fourier--Motzkin elimination without choice \\
\texttt{Uniqueness} & 100 & uniqueness of the calculus \\
\texttt{LinearLogicObstruction} & 78 & tensor collapse \\
\texttt{PCSMixture} & 345 & first order, data types, $\mathbf{n} \multimap \mathbf{m}$ \\
\texttt{Coherence} & 184 & certain states and coherence spaces \\
\texttt{Exponential} & 112 & $!\mathrm{Bool}$ \\
\texttt{Linear} & 384 & formulas, their PCS and coherence \\
\texttt{LinearInstances} & 186 & concrete results as statements about formulas \\
\texttt{Sequential} & 395 & sequential programs and the pentagon bound \\
\texttt{Kuhn}, \texttt{GeneralKuhn} & 721 & Kuhn's theorem at the second-order type \\
\texttt{General} & 389 & second order over arbitrary finite data types \\
\texttt{ProgLang} & 235 & a language for $T$ \\
\texttt{PCSGraph}, \texttt{PerfectGraph} & 314 & $\STAB$, $\QSTAB$, odd holes and antiholes \\
\texttt{ThreeArg} & 237 & a gap among total programs \\
\texttt{Retract} & 501 & coordinate retracts and the gap at other formulas \\
\texttt{Closure} & 284 & closure of the absence of a gap under $\otimes$ and $\mathbin{\&}$ \\
\texttt{CographBasic}, \texttt{CographCore} & 1439 & forests, cographs, the decomposition \\
\texttt{CographPCS}, \texttt{CographClass} & 629 & the exact condition at second order \\
\texttt{ThirdOrder} & 160 & a gap without two four-cycles \\
\texttt{TensorGap} & 81 & Shannon's clique outside the PCS \\
\texttt{LocalCalculus} & 324 & local calculi in linear and quantum logic \\
\texttt{Categorical} & 461 & duality and tightness (in part) \\
\texttt{Bell} & 214 & the Popescu--Rohrlich box at the core \\
\texttt{Quantum} & 203 & Cabello's Kochen--Specker set \\
\bottomrule
\end{tabular}
\end{center}

The general theorems are proved for arbitrary finite data types, and
for formulas built from the unit data type with $\mathbin{\&}$ and
$\oplus$ in the exact condition at second order. The finite facts behind
the witnesses are checked by evaluation with \lstinline{decide} on
explicit finite types: the five-cycle of the pentagon point, the absence
of three pairwise locally orthogonal events in the support of the
Popescu--Rohrlich box, the $64$ triples of total functions for the gap
among total programs, and the bases of the Kochen--Specker set. The
proof of the exact condition at second order uses classical case
distinctions on real numbers through the library.

\subsection{Calculi from Certain States}

\phantomsection\label{lean:mix}
\begin{lstlisting}
def IsSharpVec (x : E → ℝ≥0∞) : Prop := ∀ e, x e = 0 ∨ x e = 1

def Mix (S : Set (E → ℝ≥0∞)) : Set (E → ℝ≥0∞) :=
  {x | ∃ (n : ℕ) (w : Fin n → ℝ≥0∞) (s : Fin n → E → ℝ≥0∞),
    (∀ i, s i ∈ S) ∧ ∑ i, w i = 1 ∧ x = ∑ i, w i • s i}
\end{lstlisting}

\phantomsection\label{lean:sharp-mem-mix-iff}
\begin{lstlisting}
theorem sharp_mem_Mix_iff {S : Set (E → ℝ≥0∞)} (hS : ∀ s ∈ S, IsSharpVec s)
    {x : E → ℝ≥0∞} (hx : IsSharpVec x) : x ∈ Mix S ↔ x ∈ S
\end{lstlisting}

\phantomsection\label{lean:mix-mix}
\begin{lstlisting}
theorem Mix_Mix (S : Set (E → ℝ≥0∞)) : Mix (Mix S) = Mix S
\end{lstlisting}

\phantomsection\label{lean:affinelaw}
\begin{lstlisting}


structure AffineLaw (E : Type*) [Fintype E] where
  a : E → ℝ≥0∞
  b : E → ℝ≥0∞
  k : ℝ≥0∞
  k' : ℝ≥0∞

def AffineLaw.Holds [Fintype E] (L : AffineLaw E) (x : E → ℝ≥0∞) : Prop :=
  ∑ e, L.a e * x e + L.k ≤ ∑ e, L.b e * x e + L.k'
\end{lstlisting}

\phantomsection\label{lean:affinelaw-holds-of-mix}
\begin{lstlisting}
theorem AffineLaw.holds_of_Mix [Fintype E] (L : AffineLaw E) {S : Set (E → ℝ≥0∞)}
    (hL : ∀ s ∈ S, L.Holds s) {x : E → ℝ≥0∞} (hx : x ∈ Mix S) : L.Holds x
\end{lstlisting}

\phantomsection\label{lean:mix-eq-laws}
\begin{lstlisting}
theorem Mix_eq_laws' {S : Set (E → ℝ≥0∞)} (hS : ∀ s ∈ S, IsSharpVec s) :
    Mix S = {x | ∀ L : AffineLaw E, (∀ s ∈ S, L.Holds s) → L.Holds x}
\end{lstlisting}

\phantomsection\label{lean:frame-recipe}
\begin{lstlisting}
theorem frame_recipe :
    Set.range (fun v : Valuation (LowerSet P) => v.toFun)
      = Mix {f | ∃ v : Valuation (LowerSet P), v.IsSharp ∧ f = v.toFun}
\end{lstlisting}

\phantomsection\label{lean:constr-farkas}
\begin{lstlisting}
def Sat {n : Nat} (r : Row n) (w : Fin n → Int) : Prop :=
  0 ≤ dot n r.c w ∧ (r.s = true → 0 < dot n r.c w)

def comb {n : Nat} (p q : Int) (r₁ r₂ : Row n) : Row n :=
  ⟨fun i => p * r₁.c i + q * r₂.c i, (decide (0 < p) && r₁.s) || (decide (0 < q) && r₂.s)⟩

def zeroS (n : Nat) : Row n := ⟨fun _ => 0, true⟩

def lastc (r : Row (n + 1)) : Int := r.c (Fin.last n)

def proj (r : Row (n + 1)) : Row n := ⟨fun i => r.c i.castSucc, r.s⟩

def posR (rows : List (Row (n + 1))) := rows.filter fun r => decide (0 < lastc r)
def negR (rows : List (Row (n + 1))) := rows.filter fun r => decide (lastc r < 0)
def zerR (rows : List (Row (n + 1))) := rows.filter fun r => decide (lastc r = 0)

def pair (p q : Row (n + 1)) : Row n := proj (comb (-lastc q) (lastc p) p q)

def step (rows : List (Row (n + 1))) : List (Row n) :=
  (zerR rows).map proj ++ (posR rows).flatMap fun p => (negR rows).map fun q => pair p q

def extd (D W : Int) (w' : Fin n → Int) : Fin (n + 1) → Int :=
  fun i => if h : i.val < n then D * w' ⟨i.val, h⟩ else W

def hrow (μ : Fin N → Int) (ν α : Int) (β : Fin m → Int) : Fin (N + 1) → Int :=
  fun i => if h : i.val < N then μ ⟨i.val, h⟩ + α + d * dot m β (S ⟨i.val, h⟩)
    else ν - α - dot m β y

def hullRows : List (Row (N + 1)) :=
  (List.finRange N).map (fun j => ⟨hrow S y d (fun k => if k = j then 1 else 0) 0 0 (fun _ => 0), false⟩) ++
  [⟨hrow S y d (fun _ => 0) 1 0 (fun _ => 0), true⟩,
   ⟨hrow S y d (fun _ => 0) 0 1 (fun _ => 0), false⟩,
   ⟨hrow S y d (fun _ => 0) 0 (-1) (fun _ => 0), false⟩] ++
  (List.finRange m).flatMap (fun e =>
    [⟨hrow S y d (fun _ => 0) 0 0 (fun k => if k = e then 1 else 0), false⟩,
     ⟨hrow S y d (fun _ => 0) 0 0 (fun k => if k = e then -1 else 0), false⟩])

structure Row (n : Nat) where
  c : Fin n → Int
  s : Bool

inductive Der {n : Nat} (rows : List (Row n)) : Row n → Prop
  | base {r : Row n} : r ∈ rows → Der rows r
  | comb {r₁ r₂ : Row n} (p q : Int) : 0 ≤ p → 0 ≤ q → Der rows r₁ → Der rows r₂ →
      Der rows (comb p q r₁ r₂)

theorem fm : ∀ (n : Nat) (rows : List (Row n)), (∃ w, ∀ r ∈ rows, Sat r w) ∨ Der rows (zeroS n)

theorem complete (hd : 0 < d)
    (hlaw : ∀ (c : Fin m → Int) (k : Int), (∀ j, k ≤ dot m c (S j)) → k * d ≤ dot m c y) :
    ∃ (lam : Fin N → Int) (t : Int), 0 < t ∧ (∀ j, 0 ≤ lam j) ∧
      dot N (fun _ => 1) lam = t ∧ ∀ e, d * dot N lam (fun j => S j e) = t * y e

theorem sound (hd : 0 < d) (lam : Fin N → Int) (t : Int) (ht : 0 < t) (hl : ∀ j, 0 ≤ lam j)
    (hsum : dot N (fun _ => 1) lam = t) (hx : ∀ e, d * dot N lam (fun j => S j e) = t * y e)
    (c : Fin m → Int) (k : Int) (hc : ∀ j, k ≤ dot m c (S j)) : k * d ≤ dot m c y
\end{lstlisting}

\phantomsection\label{lean:calculus-unique}
\begin{lstlisting}
theorem calculus_unique {S C : Set (E → ℝ≥0∞)} (hS : ∀ s ∈ S, IsSharpVec s) (hSC : S ⊆ C)
    (hconv : Mix C ⊆ C)
    (hlaws : ∀ L : AffineLaw E, (∀ s ∈ S, L.Holds s) → ∀ x ∈ C, L.Holds x) :
    C = Mix S
\end{lstlisting}

\phantomsection\label{lean:calculus-laws-iff}
\begin{lstlisting}
theorem calculus_laws_iff {S C : Set (E → ℝ≥0∞)} (hSC : S ⊆ C)
    (hlaws : ∀ L : AffineLaw E, (∀ s ∈ S, L.Holds s) → ∀ x ∈ C, L.Holds x) (L : AffineLaw E) :
    (∀ x ∈ C, L.Holds x) ↔ ∀ s ∈ S, L.Holds s
\end{lstlisting}

\phantomsection\label{lean:tensor-valuation}
\begin{lstlisting}
def IsTensorValuation (v : Ω → ℝ≥0∞) (t : Ω → Ω → Ω) : Prop :=
  ∀ a b, v a + v b = v (a ⊔ b) + v (t a b)
\end{lstlisting}

\noindent Theorem~\ref{thm:tensor-collapse}.

\phantomsection\label{lean:tensor-collapse}
\begin{lstlisting}
theorem IsTensorValuation.collapse {v : Ω → ℝ≥0∞} {t : Ω → Ω → Ω}
    (h : IsTensorValuation v t) {a : Ω} (ha : v a ≠ ⊤) :
    v (t a a) = v a
\end{lstlisting}

\noindent Example~\ref{ex:tensor-witness}.

\phantomsection\label{lean:tensor-witness}
\begin{lstlisting}
theorem forced_constant_of_add_valuation {v : ℕ → ℝ≥0∞}
    (h : IsTensorValuation v (· + ·)) (hfin : ∀ n, v n ≠ ⊤) (n : ℕ) :
    v n = v 0
\end{lstlisting}

\noindent Theorem~\ref{thm:tensor-collapse}, second form.

\phantomsection\label{lean:tensor-dichotomy}
\begin{lstlisting}
theorem not_isTensorValuation_of_discriminates {v : Ω → ℝ≥0∞}
    {t : Ω → Ω → Ω} {a : Ω} (ha : v a ≠ ⊤) (hne : v (t a a) ≠ v a) :
    ¬ IsTensorValuation v t
\end{lstlisting}

\subsection{Linear Logic}

\phantomsection\label{lean:mem-lin-iff}
\begin{lstlisting}
def flat (X : Type*) [Fintype X] : Set (X → ℝ≥0∞) := {x | ∑ a, x a ≤ 1}

noncomputable def app (t : X × Y → ℝ≥0∞) (x : X → ℝ≥0∞) : Y → ℝ≥0∞ :=
  fun b => ∑ a, x a * t (a, b)

def lin (X Y : Type*) [Fintype X] [Fintype Y] : Set (X × Y → ℝ≥0∞) :=
  {t | ∀ x ∈ flat X, app t x ∈ flat Y}

theorem mem_lin_iff [DecidableEq X] (t : X × Y → ℝ≥0∞) :
    t ∈ lin X Y ↔ ∀ a, ∑ b, t (a, b) ≤ 1
\end{lstlisting}

\phantomsection\label{lean:ispcs-lin}
\begin{lstlisting}
noncomputable def pairing (x y : E → ℝ≥0∞) : ℝ≥0∞ := ∑ e, x e * y e

def row (a : X) : X × Y → ℝ≥0∞ := fun ab => if ab.1 = a then 1 else 0

def orth (S : Set (E → ℝ≥0∞)) : Set (E → ℝ≥0∞) := {y | ∀ x ∈ S, pairing x y ≤ 1}

structure IsPCS (P : Set (E → ℝ≥0∞)) : Prop where
  closed : orth (orth P) = P
  total : ∀ e, ∃ x ∈ P, x e ≠ 0
  bounded : ∀ e, ∃ M ≠ ⊤, ∀ x ∈ P, x e ≤ M

theorem isPCS_lin : IsPCS (lin X Y)

theorem isPCS_flat : IsPCS (flat X)

theorem isPCS_tensDual : IsPCS tensDual
\end{lstlisting}

\phantomsection\label{lean:pcs-recipe}
\begin{lstlisting}
theorem pcs_recipe {X Y : Type} [Fintype X] [Fintype Y] [DecidableEq X] :
    lin X Y = Mix (Set.range (det : (X → Option Y) → X × Y → ℝ≥0∞))
\end{lstlisting}

\phantomsection\label{lean:sharp-lin-iff-clique}
\begin{lstlisting}
def flatC (X : Type*) : X → X → Prop := fun x y => x = y

def dualC {X : Type*} (c : X → X → Prop) : X → X → Prop := fun x y => x = y ∨ ¬ c x y

def tensC {X Y : Type*} (c : X → X → Prop) (d : Y → Y → Prop) : X × Y → X × Y → Prop :=
  fun p q => c p.1 q.1 ∧ d p.2 q.2

def arrowC {X Y : Type*} (c : X → X → Prop) (d : Y → Y → Prop) : X × Y → X × Y → Prop :=
  dualC (tensC c (dualC d))

def IsClique {X : Type*} (c : X → X → Prop) (B : X → Prop) : Prop := ∀ x y, B x → B y → c x y

def secondC : W × W → W × W → Prop :=
  dualC (tensC (arrowC (flatC Bool) (flatC Bool)) (arrowC (flatC Bool) (flatC Bool)))

noncomputable def throughAll (C : Finset W) : Bool → Option Bool := fun a =>
  if h : ∃ p ∈ C, p.1 = a then some (Classical.choose h).2 else none

theorem sharp_lin_iff_clique {X Y : Type*} [Fintype X] [Fintype Y] [DecidableEq X] [DecidableEq Y]
    {t : X × Y → ℝ≥0∞} (h01 : ∀ p, t p = 0 ∨ t p = 1) :
    t ∈ lin X Y ↔ IsClique (arrowC (flatC X) (flatC Y)) (t · = 1)
\end{lstlisting}

\phantomsection\label{lean:linear-fm}
\begin{lstlisting}
inductive Fm where
  | base (n : ℕ)
  | neg (A : Fm)
  | tens (A B : Fm)
  | wth (A B : Fm)

def web : Fm → Type
  | base n => Fin n
  | neg A => web A
  | tens A B => web A × web B
  | wth A B => web A ⊕ web B

def P : (A : Fm) → Set (web A → ℝ≥0∞)
  | base n => flat (Fin n)
  | neg A => orth (P A)
  | tens A B => orth (orth {z | ∃ x ∈ P A, ∃ y ∈ P B, z = tvec x y})
  | wth A B => {z | (fun a => z (Sum.inl a)) ∈ P A ∧ (fun b => z (Sum.inr b)) ∈ P B}

def coh : (A : Fm) → web A → web A → Prop
  | base _ => fun x y => x = y
  | neg A => fun x y => x = y ∨ ¬ coh A x y
  | tens A B => fun s t => coh A s.1 t.1 ∧ coh B s.2 t.2
  | wth A B => fun s t => match s, t with
    | Sum.inl a, Sum.inl a' => coh A a a'
    | Sum.inr b, Sum.inr b' => coh B b b'
    | _, _ => True
\end{lstlisting}

\phantomsection\label{lean:ispcs-p}
\begin{lstlisting}
theorem isPCS_P : ∀ A : Fm, IsPCS (P A)
\end{lstlisting}

\phantomsection\label{lean:linear-invariants}
\begin{lstlisting}
def par (A B : Fm) : Fm := neg (tens (neg A) (neg B))
def plus (A B : Fm) : Fm := neg (wth (neg A) (neg B))
def lolli (A B : Fm) : Fm := neg (tens A (neg B))

noncomputable def tvec {X Y : Type} (x : X → ℝ≥0∞) (y : Y → ℝ≥0∞) : X × Y → ℝ≥0∞ :=
  fun ab => x ab.1 * y ab.2

noncomputable def svec {X Y : Type} (x : X → ℝ≥0∞) (y : Y → ℝ≥0∞) : X ⊕ Y → ℝ≥0∞
  | Sum.inl a => x a
  | Sum.inr b => y b

noncomputable def pr (s t : E) : E → ℝ≥0∞ := Pi.single s 1 + Pi.single t 1

theorem invariants : ∀ A : Fm,
    (∀ s : web A, Pi.single s 1 ∈ P A ∧ ∀ x ∈ P A, x s ≤ 1) ∧
    (∀ s t : web A, s ≠ t → coh A s t → pr s t ∈ P A) ∧
    (∀ s t : web A, s ≠ t → ¬ coh A s t → ∀ x ∈ P A, x s + x t ≤ 1)

theorem single_mem_P (A : Fm) (s : web A) : Pi.single s 1 ∈ P A

theorem pair_mem_P (A : Fm) {s t : web A} (hst : s ≠ t) (h : coh A s t) : pr s t ∈ P A
\end{lstlisting}

\phantomsection\label{lean:sharp-isclique}
\begin{lstlisting}
theorem sharp_isClique (A : Fm) {z : web A → ℝ≥0∞} (hz : z ∈ P A)
    (h01 : ∀ s, z s = 0 ∨ z s = 1) : ∀ s t, z s = 1 → z t = 1 → coh A s t
\end{lstlisting}

\phantomsection\label{lean:recipe-subset-p}
\begin{lstlisting}
theorem recipe_subset_P (A : Fm) :
    Mix {z | z ∈ P A ∧ ∀ s, z s = 0 ∨ z s = 1} ⊆ P A

theorem recipe_sharp_iff (A : Fm) {x : web A → ℝ≥0∞} (hx : ∀ s, x s = 0 ∨ x s = 1) :
    x ∈ Mix {z | z ∈ P A ∧ ∀ s, z s = 0 ∨ z s = 1} ↔ x ∈ P A
\end{lstlisting}

\phantomsection\label{lean:mix-subset-orth}
\begin{lstlisting}
theorem Mix_subset_orth {S T : Set (E → ℝ≥0∞)} (h : S ⊆ orth T) : Mix S ⊆ orth T
\end{lstlisting}

\phantomsection\label{lean:p-lolli-base}
\begin{lstlisting}
theorem P_lolli_base (n m : ℕ) : P (lolli (base n) (base m)) = lin (Fin n) (Fin m)
\end{lstlisting}

\phantomsection\label{lean:recipe-formula}
\begin{lstlisting}
def boolFin {k : ℕ} (hk : 2 ≤ k) : Bool → Fin k

theorem recipe_first_order (n m : ℕ) :
    P (lolli (base n) (base m)) = ConstructiveProb.Mix (Set.range
      (det : (Fin n → Option (Fin m)) → Fin n × Fin m → ℝ≥0∞))

theorem recipe_second_order {n m n' m' : ℕ} (hn : 0 < n) (hn' : 0 < n') :
    ConstructiveProb.Mix {z | z ∈ P (second n m n' m') ∧ ∀ u, z u = 0 ∨ z u = 1}
      = sequentialG (Fin n) (Fin m) (Fin n') (Fin m')

theorem recipe_second_order_ssubset {n m n' m' : ℕ} (hn : 2 ≤ n) (hm : 2 ≤ m) (hn' : 2 ≤ n')
    (hm' : 2 ≤ m') :
    ConstructiveProb.Mix {z | z ∈ P (second n m n' m') ∧ ∀ u, z u = 0 ∨ z u = 1}
      ⊂ P (second n m n' m')
\end{lstlisting}

\phantomsection\label{lean:sharp-bang-mixed-zero}
\begin{lstlisting}
theorem sharp_bang_mixed_zero {z : ℕ × ℕ → ℝ≥0∞} (hz : z ∈ bang) (h01 : ∀ m, z m = 0 ∨ z m = 1)
    {i j : ℕ} (hi : 1 ≤ i) (hj : 1 ≤ j) : z (i, j) = 0
\end{lstlisting}

\phantomsection\label{lean:prom-half-not-mix}
\begin{lstlisting}
def orthI {E : Type*} (S : Set (E → ℝ≥0∞)) : Set (E → ℝ≥0∞) := {y | ∀ x ∈ S, ∑' e, x e * y e ≤ 1}

def gens : Set (ℕ × ℕ → ℝ≥0∞) := {u | ∃ x ∈ flat Bool, u = prom x}

noncomputable def half : Bool → ℝ≥0∞ := fun _ => 2⁻¹

noncomputable def prom (x : Bool → ℝ≥0∞) : ℕ × ℕ → ℝ≥0∞ := fun m => x true ^ m.1 * x false ^ m.2

def bang : Set (ℕ × ℕ → ℝ≥0∞) := orthI (orthI gens)

theorem prom_half_not_mix : prom half ∈ bang ∧
    prom half ∉ Mix {z | z ∈ bang ∧ ∀ m, z m = 0 ∨ z m = 1}
\end{lstlisting}

\subsection{Second Order}

\phantomsection\label{lean:sharp-td-iff}
\begin{lstlisting}
def cov (f : X → Option Y) (g : X' → Option Y') (pq : (X × Y) × (X' × Y')) : Prop :=
  f pq.1.1 = some pq.1.2 ∧ g pq.2.1 = some pq.2.2

def cp {X Y : Type} (p q : X × Y) : Prop := p = q ∨ p.1 ≠ q.1

def cp2 (s t : (X × Y) × (X' × Y')) : Prop := cp s.1 t.1 ∧ cp s.2 t.2

noncomputable def thru {X Y : Type} [DecidableEq X] (C : Finset (X × Y)) : X → Option Y :=
  fun a => if h : ∃ p ∈ C, p.1 = a then some (Classical.choose h).2 else none

def GTree (c : X') (af : Y' → X) (pq : (X × Y) × (X' × Y')) : Prop :=
  pq.2.1 = c ∧ pq.1.1 = af pq.2.2

def tD (X Y X' Y' : Type) [Fintype X] [Fintype Y] [Fintype X'] [Fintype Y'] :
    Set ((X × Y) × (X' × Y') → ℝ≥0∞) :=
  {t | ∀ y ∈ lin X Y, ∀ z ∈ lin X' Y', ∑ p, ∑ q, y p * t (p, q) * z q ≤ 1}

def FTree (a : X) (cf : Y → X') (pq : (X × Y) × (X' × Y')) : Prop :=
  pq.1.1 = a ∧ pq.2.1 = cf pq.1.2

theorem sharp_tD_iff [Nonempty X] [Nonempty X'] {z : (X × Y) × (X' × Y') → ℝ≥0∞}
    (h01 : ∀ pq, z pq = 0 ∨ z pq = 1) :
    z ∈ tD X Y X' Y' ↔
      (∃ a cf, ∀ pq, z pq = 1 → FTree a cf pq) ∨ (∃ c af, ∀ pq, z pq = 1 → GTree c af pq)
\end{lstlisting}

\phantomsection\label{lean:sharp-mem-sequentialg}
\begin{lstlisting}
noncomputable def Core.val (K : Core A B C D) (a : A) (b : B) (c : C) (d : D) : ℝ≥0∞ :=
  K.μ a * K.κ a b c * K.ρ a b c d

noncomputable def Core.mass (K : Core A B C D) : ℝ≥0∞ := ∑ a, K.μ a

theorem sharp_mem_sequentialG [Nonempty X] [Nonempty X'] {z : (X × Y) × (X' × Y') → ℝ≥0∞}
    (hz : z ∈ tD X Y X' Y') (h01 : ∀ u, z u = 0 ∨ z u = 1) : z ∈ sequentialG X Y X' Y'
\end{lstlisting}

\phantomsection\label{lean:recipe-eq-sequentialg}
\begin{lstlisting}
structure Core (A B C D : Type) [Fintype A] [Fintype C] where
  μ : A → ℝ≥0∞
  κ : A → B → C → ℝ≥0∞
  ρ : A → B → C → D → ℝ≥0∞
  hκ : ∀ a b, ∑ c, κ a b c ≤ 1
  hρ : ∀ a b c d, ρ a b c d ≤ 1

def sequentialG (X Y X' Y' : Type) [Fintype X] [Fintype X'] : Set ((X × Y) × (X' × Y') → ℝ≥0∞) :=
  {t | ∃ (F : Core X Y X' Y') (G : Core X' Y' X Y), F.mass + G.mass ≤ 1 ∧
    t = fun pq => F.val pq.1.1 pq.1.2 pq.2.1 pq.2.2 + G.val pq.2.1 pq.2.2 pq.1.1 pq.1.2}

theorem recipe_eq_sequentialG [Nonempty X] [Nonempty X'] :
    Mix {z | z ∈ tD X Y X' Y' ∧ ∀ u, z u = 0 ∨ z u = 1} = sequentialG X Y X' Y'
\end{lstlisting}

\phantomsection\label{lean:sequentialg-subset-mix}
\begin{lstlisting}
def sharpSeqG (X Y X' Y' : Type) [Fintype X] [Fintype X'] : Set ((X × Y) × (X' × Y') → ℝ≥0∞) :=
  {z | z ∈ sequentialG X Y X' Y' ∧ ∀ u, z u = 0 ∨ z u = 1}

theorem sequentialG_subset_Mix : sequentialG X Y X' Y' ⊆ Mix (sharpSeqG X Y X' Y')
\end{lstlisting}

\phantomsection\label{lean:mix-sequentialg}
\begin{lstlisting}
noncomputable def compl' (m : A → ℝ≥0∞) : Option A → ℝ≥0∞
  | none => 1 - ∑ a, m a
  | some a => m a

abbrev ChoiceG (A B C D : Type) := (Unit → Option A) × (A × B → Option C) × (A × B × C × D → Bool)

noncomputable def detCoreG [DecidableEq A] [DecidableEq C] (σ : ChoiceG A B C D) : Core A B C D

theorem Mix_sequentialG : Mix (sequentialG X Y X' Y') = sequentialG X Y X' Y'
\end{lstlisting}

\phantomsection\label{lean:kuhn-bool}
\begin{lstlisting}
noncomputable def comp (m : Bool → ℝ≥0∞) : Option Bool → ℝ≥0∞
  | none => 1 - (m true + m false)
  | some b => m b

noncomputable def bern (r : ℝ≥0∞) : Bool → ℝ≥0∞
  | true => r
  | false => 1 - r

abbrev Choice := (Unit → Option Bool) × (Bool × Bool → Option Bool) × (Bool × Bool × Bool × Bool → Bool)

noncomputable def detCore (σ : Choice) (a b c d : Bool) : ℝ≥0∞ :=
  (if σ.1 () = some a then 1 else 0) * (if σ.2.1 (a, b) = some c then 1 else 0) *
    (if σ.2.2 (a, b, c, d) = true then 1 else 0)

noncomputable def detF (σ : Choice) : FFirst where
  μ a

noncomputable def detG (σ : Choice) : GFirst where
  μ c

def sharpSeq : Set (W × W → ℝ≥0∞) := {z | z ∈ sequential ∧ ∀ pq, z pq = 0 ∨ z pq = 1}

structure Beh where
  μ : Bool → ℝ≥0∞
  κ : Bool → Bool → Bool → ℝ≥0∞
  ρ : Bool → Bool → Bool → Bool → ℝ≥0∞
  hκ : ∀ a b, κ a b true + κ a b false ≤ 1
  hρ : ∀ a b c d, ρ a b c d ≤ 1

noncomputable def Beh.val (B : Beh) (a b c d : Bool) : ℝ≥0∞ := B.μ a * B.κ a b c * B.ρ a b c d

noncomputable def Beh.mass (B : Beh) : ℝ≥0∞ := B.μ true + B.μ false

def seqU : Set (W × W → ℝ≥0∞) :=
  {t | ∃ BF BG : Beh, BF.mass + BG.mass ≤ 1 ∧
    t = fun pq => BF.val pq.1.1 pq.1.2 pq.2.1 pq.2.2 + BG.val pq.2.1 pq.2.2 pq.1.1 pq.1.2}

theorem FFirst.den_mem_Mix (F : FFirst) : F.den ∈ Mix sharpSeq

theorem sequential_subset_Mix : sequential ⊆ Mix sharpSeq

theorem Mix_sequential : Mix sequential = sequential

theorem recipe_second_order_eq_sequential' :
    Mix {z | z ∈ tensDual ∧ ∀ pq, z pq = 0 ∨ z pq = 1} = sequential
\end{lstlisting}

\phantomsection\label{lean:proglang}
\begin{lstlisting}
inductive PEnd
  | halt
  | diverge
  | choose (k : Nat) (w : Fin k → ℝ≥0∞) (hw : ∑ i, w i ≤ 1) (p : Fin k → PEnd)

inductive POne (C D : Type)
  | diverge
  | choose (k : Nat) (w : Fin k → ℝ≥0∞) (hw : ∑ i, w i ≤ 1) (p : Fin k → POne C D)
  | call (c : C) (k : D → PEnd)

inductive P0 (X Y X' Y' : Type)
  | diverge
  | choose (k : Nat) (w : Fin k → ℝ≥0∞) (hw : ∑ i, w i ≤ 1) (p : Fin k → P0 X Y X' Y')
  | callF (a : X) (k : Y → POne X' Y')
  | callG (c : X') (k : Y' → POne X Y)

noncomputable def PEnd.den : PEnd → ℝ≥0∞
  | .halt => 1
  | .diverge => 0
  | .choose _ w _ p => ∑ i, w i * (p i).den

noncomputable def POne.den {C D : Type} [DecidableEq C] : POne C D → C × D → ℝ≥0∞
  | .diverge => fun _ => 0
  | .choose _ w _ p => fun cd => ∑ i, w i * (p i).den cd
  | .call c k => fun cd => if cd.1 = c then (k cd.2).den else 0

noncomputable def P0.den : P0 X Y X' Y' → (X × Y) × (X' × Y') → ℝ≥0∞
  | .diverge => fun _ => 0
  | .choose _ w _ p => fun u => ∑ i, w i * (p i).den u
  | .callF a k => fun u => if u.1.1 = a then (k u.1.2).den u.2 else 0
  | .callG c k => fun u => if u.2.1 = c then (k u.2.2).den u.1 else 0
\end{lstlisting}

\phantomsection\label{lean:proglang-sound}
\begin{lstlisting}
theorem P0.sound : ∀ p : P0 X Y X' Y', p.den ∈ sequentialG X Y X' Y'

theorem P0.den_mem_recipe [Nonempty X] [Nonempty X'] (p : P0 X Y X' Y') :
    p.den ∈ Mix {z | z ∈ tD X Y X' Y' ∧ ∀ u, z u = 0 ∨ z u = 1}
\end{lstlisting}

\phantomsection\label{lean:proglang-definable}
\begin{lstlisting}
def PEnd.ofBool (b : Bool) : PEnd := if b then .halt else .diverge

theorem P0.sharp_definable [Nonempty X] [Nonempty X'] {z : (X × Y) × (X' × Y') → ℝ≥0∞}
    (hz : z ∈ tD X Y X' Y') (h01 : ∀ u, z u = 0 ∨ z u = 1) : ∃ p : P0 X Y X' Y', p.den = z
\end{lstlisting}

\phantomsection\label{lean:proglang-eq}
\begin{lstlisting}
theorem recipe_eq_denotations [Nonempty X] [Nonempty X'] :
    Mix {z | z ∈ tD X Y X' Y' ∧ ∀ u, z u = 0 ∨ z u = 1} = Set.range (P0.den (X := X) (Y := Y)
      (X' := X') (Y' := Y'))
\end{lstlisting}

\phantomsection\label{lean:pente-not-mix}
\begin{lstlisting}
def emb (s : W × W) : (X × Y) × (X' × Y') :=
  ((ex s.1.1, ey s.1.2), (ex' s.2.1, ey' s.2.2))

noncomputable def pentE (ex : Bool → X) (ey : Bool → Y) (ex' : Bool → X') (ey' : Bool → Y') :
    (X × Y) × (X' × Y') → ℝ≥0∞ :=
  fun u => if u ∈ pent.image (emb ex ey ex' ey') then 2⁻¹ else 0

theorem pentE_not_mix (hx : Function.Injective ex) (hy : Function.Injective ey)
    (hx' : Function.Injective ex') (hy' : Function.Injective ey')
    {n : ℕ} (w : Fin n → ℝ≥0∞) (hw : ∑ i, w i = 1) (z : Fin n → (X × Y) × (X' × Y') → ℝ≥0∞)
    (hz : ∀ i, z i ∈ tD X Y X' Y') (h01 : ∀ i u, z i u = 0 ∨ z i u = 1) :
    pentE ex ey ex' ey' ≠ ∑ i, w i • z i

theorem pentE_mem (hx : Function.Injective ex) (hy : Function.Injective ey)
    (hx' : Function.Injective ex') (hy' : Function.Injective ey') :
    pentE ex ey ex' ey' ∈ tD X Y X' Y'
\end{lstlisting}

\phantomsection\label{lean:pent-cycle}
\begin{lstlisting}
noncomputable def idBool : Bool × Bool → ℝ≥0∞ := fun ab => if ab.1 = ab.2 then 1 else 0

noncomputable def det (f : X → Option Y) : X × Y → ℝ≥0∞

noncomputable def rowDist (t : X × Y → ℝ≥0∞) (a : X) : Option Y → ℝ≥0∞
  | none => 1 - ∑ b, t (a, b)
  | some b => t (a, b)

noncomputable def tens (x : X → ℝ≥0∞) (y : Y → ℝ≥0∞) : X × Y → ℝ≥0∞ := fun ab => x ab.1 * y ab.2

abbrev W := Bool × Bool

def pent : Finset (W × W) :=
  {((true, true), (true, true)), ((true, false), (false, false)),
   ((false, true), (true, false)), ((false, true), (false, true)),
   ((false, false), (false, false))}

noncomputable def pentElt : W × W → ℝ≥0∞ := fun pq => if pq ∈ pent then 2⁻¹ else 0

def covers (f g : Bool → Option Bool) (pq : W × W) : Prop :=
  f pq.1.1 = some pq.1.2 ∧ g pq.2.1 = some pq.2.2

def compat (p q : W) : Prop := p = q ∨ p.1 ≠ q.1

def compat2 (s t : W × W) : Prop := compat s.1 t.1 ∧ compat s.2 t.2

def through (p q : W) : Bool → Option Bool := fun b => if b = p.1 then some p.2 else some q.2

theorem pent_clique_le_two : ∀ B ∈ pent.powerset,
    (∀ s ∈ B, ∀ t ∈ B, compat2 s t) → B.card ≤ 2

theorem pent_indep_le_two : ∀ B ∈ pent.powerset,
    (∀ s ∈ B, ∀ t ∈ B, compat2 s t → s = t) → B.card ≤ 2
\end{lstlisting}

\phantomsection\label{lean:pent-bounds}
\begin{lstlisting}
structure FFirst where
  μ : Bool → ℝ≥0∞
  κ : Bool → Bool → Bool → ℝ≥0∞
  ρ : Bool → Bool → Bool → Bool → ℝ≥0∞
  hμ : μ true + μ false ≤ 1
  hκ : ∀ a b, κ a b true + κ a b false ≤ 1
  hρ : ∀ a b c d, ρ a b c d ≤ 1

structure GFirst where
  μ : Bool → ℝ≥0∞
  κ : Bool → Bool → Bool → ℝ≥0∞
  ρ : Bool → Bool → Bool → Bool → ℝ≥0∞
  hμ : μ true + μ false ≤ 1
  hκ : ∀ c d, κ c d true + κ c d false ≤ 1
  hρ : ∀ c d a b, ρ c d a b ≤ 1

noncomputable def FFirst.den (S : FFirst) : W × W → ℝ≥0∞ := fun pq =>
  S.μ pq.1.1 * S.κ pq.1.1 pq.1.2 pq.2.1 * S.ρ pq.1.1 pq.1.2 pq.2.1 pq.2.2

noncomputable def GFirst.den (S : GFirst) : W × W → ℝ≥0∞ := fun pq =>
  S.μ pq.2.1 * S.κ pq.2.1 pq.2.2 pq.1.1 * S.ρ pq.2.1 pq.2.2 pq.1.1 pq.1.2

noncomputable def copyStrategy : FFirst where
  μ a

def idleG : GFirst where
  μ _

def idleF : FFirst where
  μ _

def tensDual : Set (W × W → ℝ≥0∞) :=
  {t | ∀ y ∈ lin Bool Bool, ∀ z ∈ lin Bool Bool, ∑ p, ∑ q, y p * t (p, q) * z q ≤ 1}

def sequential : Set (W × W → ℝ≥0∞) :=
  {t | ∃ (α β : ℝ≥0∞) (F : FFirst) (G : GFirst), α + β ≤ 1 ∧ t = α • F.den + β • G.den}

theorem seq_pent_le_two {t : W × W → ℝ≥0∞} (ht : t ∈ sequential) : ∑ pq ∈ pent, t pq ≤ 2

theorem seq_pent_tight : ∃ t ∈ sequential, ∑ pq ∈ pent, t pq = 2

theorem pentElt_pent_sum : ∑ pq ∈ pent, pentElt pq = 5 * 2⁻¹

theorem tensDual_pent_le {t : W × W → ℝ≥0∞} (ht : t ∈ tensDual) :
    2 * ∑ pq ∈ pent, t pq ≤ 5
\end{lstlisting}

\phantomsection\label{lean:p-second}
\begin{lstlisting}
def second (n m n' m' : ℕ) : Fm := neg (tens (lolli (base n) (base m)) (lolli (base n') (base m')))

theorem P_second (n m n' m' : ℕ) :
    P (second n m n' m') = tD (Fin n) (Fin m) (Fin n') (Fin m')
\end{lstlisting}

\phantomsection\label{lean:gap-t3}
\begin{lstlisting}
abbrev L2 : Fm := lolli (base 2) (base 2)

abbrev S3 : Fm := tens (tens L2 L2) L2

abbrev T3 : Fm := Fm.neg S3

noncomputable def graph2 (f : Fin 2 → Fin 2) : web L2 → ℝ≥0∞ := fun p => if f p.1 = p.2 then 1 else 0

theorem graph2_mem (f : Fin 2 → Fin 2) : graph2 f ∈ P L2

theorem sharp_L2_le {s : web L2 → ℝ≥0∞} (hs : s ∈ sharpP L2) : ∃ f, s ≤ graph2 f

theorem sharp_L2L2_le {s : web (tens L2 L2) → ℝ≥0∞} (hs : s ∈ sharpP (tens L2 L2)) :
    ∃ f g, s ≤ tvec (graph2 f) (graph2 g)

def k3 (v : W3) : ℕ := if v = onePt then 2 else if v ∈ halfPts then 1 else 0

noncomputable def y5 : web T3 → ℝ≥0∞ := fun v => (k3 v : ℝ≥0∞) * 2⁻¹

theorem y5_total (f g h : Fin 2 → Fin 2) :
    pairing (tvec (tvec (graph2 f) (graph2 g)) (graph2 h)) y5 = 1

theorem y5_mem : y5 ∈ P T3

theorem sharp5_le {s : web T3 → ℝ≥0∞} (hs : s ∈ P T3) (h01 : ∀ v, s v = 0 ∨ s v = 1) :
    s c0 + s c1 + s c2 + s c3 + s c4 ≤ 2

theorem gap_T3 : recipe T3 ⊂ P T3
\end{lstlisting}

\subsection{The Exact Condition}

\phantomsection\label{lean:td-eq-qstab}
\begin{lstlisting}
def excl (s t : (X × Y) × (X' × Y')) : Prop := s ≠ t ∧ cp2 s t

theorem tD_eq_QSTAB : tD X Y X' Y' = QSTAB (excl (X := X) (Y := Y) (X' := X') (Y' := Y'))
\end{lstlisting}

\phantomsection\label{lean:recipe-eq-stab}
\begin{lstlisting}
theorem recipe_eq_STAB :
    Mix {z | z ∈ tD X Y X' Y' ∧ ∀ u, z u = 0 ∨ z u = 1}
      = STAB (excl (X := X) (Y := Y) (X' := X') (Y' := Y'))
\end{lstlisting}

\phantomsection\label{lean:berge-of-stab-eq-qstab}
\begin{lstlisting}
def cyc {n : ℕ} (i j : ZMod n) : Prop := j = i + 1 ∨ i = j + 1

def IsCliqueG (C : Finset V) : Prop := ∀ u ∈ C, ∀ v ∈ C, u ≠ v → adj u v

def IsStableG (I : Finset V) : Prop := ∀ u ∈ I, ∀ v ∈ I, ¬ adj u v

noncomputable def ind (I : Finset V) : V → ℝ≥0∞ := fun v => if v ∈ I then 1 else 0

structure Antihole (adj : V → V → Prop) (n : ℕ) where
  h : ZMod n → V
  inj : Function.Injective h
  adj_iff : ∀ i j, i ≠ j → (adj (h i) (h j) ↔ ¬ cyc i j)

def QSTAB : Set (V → ℝ≥0∞) := {x | ∀ C, IsCliqueG adj C → ∑ v ∈ C, x v ≤ 1}

def STAB : Set (V → ℝ≥0∞) := Mix {x | ∃ I, IsStableG adj I ∧ x = ind I}

structure Hole (adj : V → V → Prop) (n : ℕ) where
  h : ZMod n → V
  inj : Function.Injective h
  adj_iff : ∀ i j, i ≠ j → (adj (h i) (h j) ↔ cyc i j)

theorem berge_of_stab_eq_qstab (heq : STAB adj = QSTAB adj) (k : ℕ) (hk : 2 ≤ k)
    [NeZero (2 * k + 1)] : IsEmpty (Hole adj (2 * k + 1)) ∧ IsEmpty (Antihole adj (2 * k + 1))
\end{lstlisting}

\phantomsection\label{lean:berge-of-recipe-eq-pcs}
\begin{lstlisting}
theorem berge_of_recipe_eq_pcs
    (heq : Mix {z | z ∈ tD X Y X' Y' ∧ ∀ u, z u = 0 ∨ z u = 1} = tD X Y X' Y')
    (k : ℕ) (hk : 2 ≤ k) [NeZero (2 * k + 1)] :
    IsEmpty (Hole (excl (X := X) (Y := Y) (X' := X') (Y' := Y')) (2 * k + 1)) ∧
      IsEmpty (Antihole (excl (X := X) (Y := Y) (X' := X') (Y' := Y')) (2 * k + 1))
\end{lstlisting}

\phantomsection\label{lean:closure}
\begin{lstlisting}
theorem orth_orth_subset_Mix {S : Set (E → ℝ≥0∞)} (hSf : S.Finite)
    (hS : ∀ s ∈ S, ConstructiveProb.IsSharpVec s) (h0 : (0 : E → ℝ≥0∞) ∈ S)
    (hdown : ∀ s ∈ S, ∀ t : E → ℝ≥0∞, ConstructiveProb.IsSharpVec t → t ≤ s → t ∈ S) :
    orth (orth S) ⊆ ConstructiveProb.Mix S

theorem orth_Mix (S : Set (E → ℝ≥0∞)) : orth (ConstructiveProb.Mix S) = orth S

theorem recipe_orth_orth (A : Fm) : orth (orth (recipe A)) = recipe A

theorem recipe_tens {A B : Fm} (hA : recipe A = P A) (hB : recipe B = P B) :
    recipe (tens A B) = P (tens A B)

theorem recipe_wth {A B : Fm} (hA : recipe A = P A) (hB : recipe B = P B) :
    recipe (wth A B) = P (wth A B)

def sharpP (A : Fm) : Set (web A → ℝ≥0∞) := {z | z ∈ P A ∧ ∀ s, z s = 0 ∨ z s = 1}

theorem P_neg_tens_eq {A B : Fm} (hA : recipe A = P A) (hB : recipe B = P B) :
    P (Fm.neg (tens A B)) = orth {z | ∃ s ∈ sharpP A, ∃ t ∈ sharpP B, z = tvec s t}

def svec' {X Y : Type} (x : X → ℝ≥0∞) (y : Y → ℝ≥0∞) : X ⊕ Y → ℝ≥0∞
  | Sum.inl a => x a
  | Sum.inr b => y b

theorem tvec_sum {X Y : Type} [Fintype X] [Fintype Y] {m n : ℕ} (w : Fin m → ℝ≥0∞)
    (s : Fin m → X → ℝ≥0∞) (v : Fin n → ℝ≥0∞) (t : Fin n → Y → ℝ≥0∞) :
    tvec (∑ i, w i • s i) (∑ j, v j • t j)
      = ∑ ij : Fin m × Fin n, (w ij.1 * v ij.2) • tvec (s ij.1) (t ij.2)
\end{lstlisting}

\phantomsection\label{lean:recipe-core}
\begin{lstlisting}
theorem recipe_lolli_base (n m : ℕ) :
    recipe (lolli (base n) (base m)) = P (lolli (base n) (base m))

theorem recipe_core (n m n' m' : ℕ) : recipe (core n m n' m') = P (core n m n' m')
\end{lstlisting}

\phantomsection\label{lean:recipe-code}
\begin{lstlisting}
inductive ACode
  | one
  | wth (a b : ACode)
  | plus (a b : ACode)

def toFm : ACode → Fm
  | one => Fm.base 1
  | wth a b => Fm.wth a.toFm b.toFm
  | plus a b => Fm.plus a.toFm b.toFm

abbrev W (a : ACode) : Type := web a.toFm

theorem recipe_plus {A B : Fm} (hA : recipe A = P A) (hB : recipe B = P B) :
    recipe (Fm.plus A B) = P (Fm.plus A B)

theorem recipe_code : ∀ a : ACode, recipe a.toFm = P a.toFm

theorem clique_mem_P : ∀ (a : ACode) (c : W a → ℝ≥0∞), (∀ x, c x = 0 ∨ c x = 1) →
    (∀ x y, c x = 1 → c y = 1 → ACode.coh a x y) → c ∈ P a.toFm
\end{lstlisting}

\phantomsection\label{lean:tree-of-noc4}
\begin{lstlisting}
inductive TCode
  | one
  | plus (s t : TCode)
  | cone (t : TCode)

def toA : TCode → ACode
  | one => ACode.one
  | plus s t => ACode.plus s.toA t.toA
  | cone t => ACode.wth t.toA ACode.one

structure GIso (a b : ACode) where
  e : W a ≃ W b
  coh_iff : ∀ x y, ACode.coh a x y ↔ ACode.coh b (e x) (e y)

def GIso.retract {a b : ACode} (I : GIso a b) : Retract a.toFm b.toFm

def retractSwap (A B : Fm) : Retract (Fm.tens A B) (Fm.tens B A)

def HasC4 (a : ACode) : Prop :=
  ∃ p q r s : W a, ACode.coh a p q ∧ ACode.coh a q r ∧ ACode.coh a r s ∧ ACode.coh a s p ∧
    ¬ ACode.coh a p r ∧ ¬ ACode.coh a q s

theorem hasC4_iff (a : ACode) : HasC4 a ↔ Nonempty (C4 a.toFm)

def Complete (a : ACode) : Prop := ∀ x y, ACode.coh a x y

theorem complete_or {a b : ACode} (h : ¬ HasC4 (ACode.wth a b)) : Complete a ∨ Complete b

theorem tree_of_noC4 : ∀ a : ACode, ¬ HasC4 a → ∃ t : TCode, Nonempty (GIso a t.toA)
\end{lstlisting}

\phantomsection\label{lean:mem-p-iff}
\begin{lstlisting}
noncomputable def ω : (a : ACode) → (W a → ℝ) → ℝ
  | one, v => v pt
  | wth a b, v => ω a (fun x => v (Sum.inl x)) + ω b (fun x => v (Sum.inr x))
  | plus a b, v => max (ω a (fun x => v (Sum.inl x))) (ω b (fun x => v (Sum.inr x)))

abbrev N (t : TCode) : Type := ACode.W t.toA

def PathSums {V : Type*} [AddCommMonoid V] : (t : TCode) → (N t → V) → Set V
  | one, y => {0, y ACode.pt}
  | plus s t, y => PathSums s (fun x => y (Sum.inl x)) ∪ PathSums t (fun x => y (Sum.inr x))
  | cone t, y => insert 0 ((fun S => y (root t) + S) '' PathSums t (fun x => y (Sum.inl x)))

def IsPathInd : (t : TCode) → (N t → ℝ) → Prop
  | TCode.one, a => a = 0 ∨ a = fun _ => 1
  | TCode.plus s u, a =>
      (∃ a1 : N s → ℝ, IsPathInd s a1 ∧ a = Sum.elim a1 (0 : N u → ℝ)) ∨
        (∃ a2 : N u → ℝ, IsPathInd u a2 ∧ a = Sum.elim (0 : N s → ℝ) a2)
  | TCode.cone u, a => a = 0 ∨
      ∃ a' : N u → ℝ, IsPathInd u a' ∧ a = Sum.elim a' (fun _ : W ACode.one => (1 : ℝ))

theorem clique_le_path : ∀ (t : TCode) (c : N t → ℝ), (∀ x, c x = 0 ∨ c x = 1) →
    (∀ x y, c x = 1 → c y = 1 → ACode.coh t.toA x y) → ∃ a, IsPathInd t a ∧ c ≤ a

theorem mem_P_iff (t : TCode) (K : ACode) (Z : N t → W K → ℝ) (hZ : ∀ x k, 0 ≤ Z x k) :
    (fun p : N t × W K => ENNReal.ofReal (Z p.1 p.2)) ∈
        P (Fm.neg (Fm.tens t.toA.toFm K.toFm)) ↔
      ∀ S ∈ PathSums t Z, ω K S ≤ 1
\end{lstlisting}

\phantomsection\label{lean:decompose}
\begin{lstlisting}
def PMix (S : Set E) : Set E :=
  {x | ∃ (ι : Type) (_ : Fintype ι) (w : ι → ℝ) (s : ι → E),
    (∀ i, 0 < w i) ∧ ∑ i, w i = 1 ∧ (∀ i, s i ∈ S) ∧ x = ∑ i, w i • s i}

def IsStab (a : ACode) (v : W a → ℝ) : Prop := (∀ x, v x = 0 ∨ v x = 1) ∧ ω a v ≤ 1

def RootSet (K : ACode) (t : TCode) : Set ((W K → ℝ) × (N t → W K → ℝ)) :=
  {p | IsStab K p.1 ∧ 0 ≤ p.2 ∧ ∀ S ∈ PathSums t p.2, ω K (p.1 + S) ≤ 1}

theorem root_step : ∀ (K : ACode) (t : TCode) (m : W K → ℝ) (y : N t → W K → ℝ),
    0 ≤ m → 0 ≤ y → (∀ S ∈ PathSums t y, ω K (m + S) ≤ 1) → (m, y) ∈ PMix (RootSet K t)

def Valid (t : TCode) (K : ACode) : Set (N t → W K → ℝ) :=
  {σ | (∀ x h, σ x h = 0 ∨ σ x h = 1) ∧ ∀ S ∈ PathSums t σ, IsStab K S}

theorem main : ∀ (t : TCode) (K : ACode) (z : N t → W K → ℝ), 0 ≤ z →
    (∀ S ∈ PathSums t z, ω K S ≤ 1) → z ∈ PMix (Valid t K)

theorem second_order_noGap (t : TCode) (K : ACode) :
    recipe (Fm.neg (Fm.tens t.toA.toFm K.toFm)) = P (Fm.neg (Fm.tens t.toA.toFm K.toFm))
\end{lstlisting}

\phantomsection\label{lean:absorb}
\begin{lstlisting}
theorem exact_absorb : ∀ (K : ACode) {b v : W K → ℝ} {L : ℝ}, 0 ≤ b → 0 ≤ v → ω K b ≤ L →
    L ≤ ω K (b + v) → ∃ a : W K → ℝ, 0 ≤ a ∧ a ≤ v ∧ ω K (b + a) = L

def AbsOK (K : ACode) (L : ℝ) (m : W K → ℝ) (S : (W K → ℝ) × (W K → ℝ)) : Prop :=
  ω K (m + S.1) ≤ L ∧ ω K S.2 + L ≤ max L (ω K (m + S.1 + S.2))

def AbsAt (K : ACode) (t : TCode) (L : ℝ) (m : W K → ℝ) (y : N t → W K → ℝ) : Prop :=
  ∃ p q : N t → W K → ℝ, 0 ≤ p ∧ 0 ≤ q ∧ p + q = y ∧
    ∀ S ∈ PathSums t (fun x => (p x, q x)), AbsOK K L m S

theorem greedy (K : ACode)
    (abs0 : ∀ (t : TCode) (m : W K → ℝ) (y : N t → W K → ℝ), 0 ≤ m → 0 ≤ y →
      AbsAt K t (ω K m) m y) :
    ∀ (t : TCode) (L : ℝ) (m : W K → ℝ) (y : N t → W K → ℝ), 0 ≤ m → 0 ≤ y → ω K m ≤ L →
      AbsAt K t L m y

theorem absorb : ∀ (K : ACode) (t : TCode) (L : ℝ) (m : W K → ℝ) (y : N t → W K → ℝ),
    0 ≤ m → 0 ≤ y → ω K m ≤ L → AbsAt K t L m y
\end{lstlisting}

\phantomsection\label{lean:second-iff}
\begin{lstlisting}
theorem second_order_iff (a b : ACode) :
    recipe (Fm.neg (Fm.tens a.toFm b.toFm)) = P (Fm.neg (Fm.tens a.toFm b.toFm)) ↔
      ¬ (HasC4 a ∧ HasC4 b)
\end{lstlisting}

\phantomsection\label{lean:gap-f3}
\begin{lstlisting}
abbrev X3 : Fm := wth (base 2) (base 1)

abbrev F3 : Fm := Fm.neg (tens X3 S2)

theorem F3_eq : F3 = lolli X3 (core 2 2 2 2) := rfl

def w0 : web S2 := ((0, 0), (0, 0))
def w1 : web S2 := ((0, 0), (0, 1))
def w2 : web S2 := ((0, 1), (1, 1))
def w3 : web S2 := ((1, 0), (0, 1))
def w4 : web S2 := ((1, 1), (1, 0))

def p0 : web X3 := Sum.inl 0
def p1 : web X3 := Sum.inl 1
def pu : web X3 := Sum.inr 0

def supp3 : Finset (web F3) := {(p0, w2), (p0, w4), (p1, w0), (pu, w1), (pu, w3)}

noncomputable def y3 : web F3 → ℝ≥0∞ := fun v => if v ∈ supp3 then 2⁻¹ else 0

theorem y3_mem : y3 ∈ P F3

theorem sharp3_le {s : web F3 → ℝ≥0∞} (hs : s ∈ P F3) (h01 : ∀ v, s v = 0 ∨ s v = 1) :
    ∑ v ∈ supp3, s v ≤ 2

theorem gap_F3 : recipe F3 ⊂ P F3

theorem X3_no_C4 (Q : C4 X3) : False
\end{lstlisting}

\subsection{Other Formulas}

\phantomsection\label{lean:retract}
\begin{lstlisting}
noncomputable def push (j : E → F) (x : E → ℝ≥0∞) : F → ℝ≥0∞ :=
  fun a => ∑ d, if j d = a then x d else 0

structure Retract (D A : Fm) where
  j : web D → web A
  inj : Function.Injective j
  push_mem : ∀ x ∈ P D, push j x ∈ P A
  pull_mem : ∀ z ∈ P A, (fun d => z (j d)) ∈ P D

def recipe (A : Fm) : Set (web A → ℝ≥0∞) :=
  ConstructiveProb.Mix {z | z ∈ P A ∧ ∀ s, z s = 0 ∨ z s = 1}
\end{lstlisting}

\phantomsection\label{lean:retract-gap}
\begin{lstlisting}
theorem Retract.recipe_ne {D A : Fm} (R : Retract D A) (h : recipe D ≠ P D) :
    recipe A ≠ P A

theorem pull_push {j : E → F} (hj : Function.Injective j) (x : E → ℝ≥0∞) :
    (fun d => push j x (j d)) = x

theorem pairing_push (j : E → F) (u : E → ℝ≥0∞) (z : F → ℝ≥0∞) :
    pairing (push j u) z = pairing u (fun d => z (j d))
\end{lstlisting}

\phantomsection\label{lean:retract-ops}
\begin{lstlisting}
def Retract.refl (D : Fm) : Retract D D

def Retract.trans {D A C : Fm} (R : Retract D A) (S : Retract A C) : Retract D C

def Retract.neg {D A : Fm} (R : Retract D A) : Retract (Fm.neg D) (Fm.neg A)

def Retract.negneg (D : Fm) : Retract D (Fm.neg (Fm.neg D))

def Retract.tensL (D B : Fm) (b : web B) : Retract D (Fm.tens D B)

def Retract.tensR (D B : Fm) (b : web B) : Retract D (Fm.tens B D)

def Retract.wthL (D B : Fm) : Retract D (Fm.wth D B)

def Retract.wthR (D B : Fm) : Retract D (Fm.wth B D)

theorem push_tensL {X Y : Type} [Fintype X] [DecidableEq X] [DecidableEq Y] (b : Y)
    (x : X → ℝ≥0∞) : push (fun d => (d, b)) x = tvec x (Pi.single b 1)

theorem push_tensR {X Y : Type} [Fintype X] [DecidableEq X] [DecidableEq Y] (b : Y)
    (x : X → ℝ≥0∞) : push (fun d => (b, d)) x = tvec (Pi.single b 1) x

theorem push_comp [Fintype F] (j1 : E → F) (j2 : F → G) (x : E → ℝ≥0∞) :
    push (fun d => j2 (j1 d)) x = push j2 (push j1 x)

theorem push_id [DecidableEq E] (x : E → ℝ≥0∞) : push (fun d : E => d) x = x
\end{lstlisting}

\phantomsection\label{lean:gap-of-occ}
\begin{lstlisting}
def NE : Fm → Prop
  | base n => 0 < n
  | Fm.neg A => NE A
  | tens A B => NE A ∧ NE B
  | wth A B => NE A ∧ NE B

theorem NE.point : ∀ A : Fm, NE A → Nonempty (web A)
  | base _, h => ⟨⟨0, h⟩⟩
  | Fm.neg A, h => NE.point A h
  | tens A B, h => by
    obtain ⟨a⟩ := NE.point A h.1
    obtain ⟨b⟩ := NE.point B h.2
    exact ⟨(a, b)⟩
  | wth A _, h => by
    obtain ⟨a⟩ := NE.point A h.1
    exact ⟨Sum.inl a⟩

/-- `Occ D p A`: `D` occurs in `A`, positively if `p` and negatively otherwise. -/
inductive Occ (D : Fm) : Bool → Fm → Prop
  | self : Occ D true D
  | neg {p : Bool} {A : Fm} : Occ D p A → Occ D (!p) (Fm.neg A)
  | tensL {p : Bool} {A : Fm} (B : Fm) : Occ D p A → Occ D p (tens A B)
  | tensR {p : Bool} {A : Fm} (B : Fm) : Occ D p A → Occ D p (tens B A)
  | wthL {p : Bool} {A : Fm} (B : Fm) : Occ D p A → Occ D p (wth A B)
  | wthR {p : Bool} {A : Fm} (B : Fm) : Occ D p A → Occ D p (wth B A)

/-- `D` or `D^⊥`, according to polarity. -/
def pol : Bool → Fm → Fm
  | true, D => D
  | false, D => Fm.neg D

/-- **Occurrences give retracts.** -/
theorem occ_retract {D : Fm} : ∀ {p : Bool} {A : Fm}, Occ D p A → NE A →
    Nonempty (Retract (pol p D) A)
  | _, _, Occ.self, _ => ⟨Retract.refl D⟩
  | _, _, Occ.neg (p := p) h, hA => by
    obtain ⟨R⟩ := occ_retract h hA
    cases p
    · exact ⟨(Retract.negneg D).trans R.neg⟩
    · exact ⟨R.neg⟩
  | _, _, Occ.tensL B h, hA => by
    obtain ⟨R⟩ := occ_retract h hA.1
    obtain ⟨b⟩ := NE.point B hA.2
    exact ⟨R.trans (Retract.tensL _ B b)⟩
  | _, _, Occ.tensR B h, hA => by
    obtain ⟨R⟩ := occ_retract h hA.2
    obtain ⟨b⟩ := NE.point B hA.1
    exact ⟨R.trans (Retract.tensR _ B b)⟩
  | _, _, Occ.wthL B h, hA => by
    obtain ⟨R⟩ := occ_retract h hA.1
    exact ⟨R.trans (Retract.wthL _ B)⟩
  | _, _, Occ.wthR B h, hA => by
    obtain ⟨R⟩ := occ_retract h hA.2
    exact ⟨R.trans (Retract.wthR _ B)⟩

/-- The second-order core `(n ⊸ m) ⊗ (n' ⊸ m')`. -/
abbrev core (n m n' m' : ℕ) : Fm := tens (lolli (base n) (base m)) (lolli (base n') (base m'))

/-- **The gap at every formula with a negative occurrence of the core.** -/
theorem gap_of_occ {n m n' m' : ℕ} (hn : 2 ≤ n) (hm : 2 ≤ m) (hn' : 2 ≤ n') (hm' : 2 ≤ m')
    {A : Fm} (hocc : Occ (core n m n' m') false A) (hA : NE A) : recipe A ⊂ P A

example : recipe (lolli (core 2 2 2 2) (base 2)) ⊂ P (lolli (core 2 2 2 2) (base 2))
\end{lstlisting}

\phantomsection\label{lean:retract-tens}
\begin{lstlisting}
theorem push_tvec {X Y X' Y' : Type} [Fintype X] [Fintype Y] [DecidableEq X'] [DecidableEq Y']
    (j : X → X') (k : Y → Y') (x : X → ℝ≥0∞) (y : Y → ℝ≥0∞) :
    push (fun d : X × Y => (j d.1, k d.2)) (tvec x y) = tvec (push j x) (push k y)

def Retract.tens {D A D' A' : Fm} (R : Retract D A) (R' : Retract D' A') :
    Retract (Fm.tens D D') (Fm.tens A A')

theorem gap_of_retracts {n m n' m' : ℕ} (hn : 2 ≤ n) (hm : 2 ≤ m) (hn' : 2 ≤ n') (hm' : 2 ≤ m')
    {X Y : Fm} (RX : Retract (lolli (base n) (base m)) X)
    (RY : Retract (lolli (base n') (base m')) Y) :
    recipe (Fm.neg (tens X Y)) ⊂ P (Fm.neg (tens X Y))
\end{lstlisting}

\phantomsection\label{lean:c4-retract}
\begin{lstlisting}
theorem coh_symm : ∀ (A : Fm) {s t : web A}, coh A s t → coh A t s

structure C4 (X : Fm) where
  (a b c d : web X)
  hab : coh X a b
  hbc : coh X b c
  hcd : coh X c d
  hda : coh X d a
  hac : ¬ coh X a c
  hbd : ¬ coh X b d

def C4.j (Q : C4 X) : Fin 2 × Fin 2 → web X := fun p => ![![Q.a, Q.c], ![Q.b, Q.d]] p.1 p.2

def C4.retract (Q : C4 X) : Retract (lolli (base 2) (base 2)) X
\end{lstlisting}

\phantomsection\label{lean:gap-of-c4}
\begin{lstlisting}
theorem gap_of_C4 {X Y : Fm} (QX : C4 X) (QY : C4 Y) :
    recipe (Fm.neg (tens X Y)) ⊂ P (Fm.neg (tens X Y))

theorem gap_of_C4_occ {X Y A : Fm} (QX : C4 X) (QY : C4 Y) (hocc : Occ (tens X Y) false A)
    (hA : NE A) : recipe A ⊂ P A
\end{lstlisting}

\phantomsection\label{lean:gap-curried}
\begin{lstlisting}
def C4.lolli22 : C4 (lolli (base 2) (base 2))

def C4.curried : C4 (Fm.neg (lolli (lolli (base 2) (base 2)) (base 2)))

theorem gap_curried :
    recipe (lolli (lolli (base 2) (base 2)) (lolli (lolli (base 2) (base 2)) (base 2)))
      ⊂ P (lolli (lolli (base 2) (base 2)) (lolli (lolli (base 2) (base 2)) (base 2)))
\end{lstlisting}

\subsection{Local Certainty and Duality}

\phantomsection\label{lean:local-calc}
\begin{lstlisting}
def locS (C : Set (Finset E)) : Set (E → ℝ≥0∞) := {x | ∀ t ∈ C, ∑ e ∈ t, x e ≤ 1}

def locN (C : Set (Finset E)) : Set (E → ℝ≥0∞) := {x | ∀ t ∈ C, ∑ e ∈ t, x e = 1}

def globS (C : Set (Finset E)) : Set (E → ℝ≥0∞) := Mix {x | x ∈ locS C ∧ IsSharpVec x}

def globN (C : Set (Finset E)) : Set (E → ℝ≥0∞) := Mix {x | x ∈ locN C ∧ IsSharpVec x}

theorem globN_sub (C : Set (Finset E)) : globN C ⊆ locN C

theorem globS_sub (C : Set (Finset E)) : globS C ⊆ locS C

theorem globN_single : globN ({Finset.univ} : Set (Finset E)) = locN {Finset.univ}

theorem globS_single : globS ({Finset.univ} : Set (Finset E)) = locS {Finset.univ}
\end{lstlisting}

\phantomsection\label{lean:loc-linear}
\begin{lstlisting}
def loc (A : Fm) : Set (web A → ℝ≥0∞) := orth (sharpP (Fm.neg A))

def dualCtx (A : Fm) : Set (Finset (web A)) :=
  {t | ∃ s ∈ sharpP (Fm.neg A), t = Finset.univ.filter fun e => s e = 1}

theorem loc_eq_locS (A : Fm) : loc A = locS (dualCtx A)

theorem recipe_sub_P_sub_loc (A : Fm) : recipe A ⊆ P A ∧ P A ⊆ loc A

theorem P_eq_loc_iff (A : Fm) : P A = loc A ↔ recipe (Fm.neg A) = P (Fm.neg A)

theorem clique_mem_loc (A : Fm) {c : web A → ℝ≥0∞} (h01 : ∀ e, c e = 0 ∨ c e = 1)
    (hc : ∀ x y, c x = 1 → c y = 1 → coh A x y) : c ∈ loc A

theorem P_S2_eq_loc : P S2 = loc S2

theorem P_TT_ssubset_loc : P TT ⊂ loc TT
\end{lstlisting}

\phantomsection\label{lean:shannon-clique}
\begin{lstlisting}
def φ : W × W → web S2 := emb (boolFin le_rfl) (boolFin le_rfl) (boolFin le_rfl) (boolFin le_rfl)

abbrev S2 : Fm := second 2 2 2 2

abbrev TT : Fm := neg (tens S2 S2)

noncomputable def shannonF : Finset (web (tens S2 S2)) := shannon.image fun u => (φ u.1, φ u.2)

theorem shannonF_clique : ∀ u ∈ shannonF, ∀ v ∈ shannonF, coh TT u v
\end{lstlisting}

\phantomsection\label{lean:shannon-not-mem}
\begin{lstlisting}
noncomputable def pentS2 : web S2 → ℝ≥0∞ :=
  pentE (boolFin le_rfl) (boolFin le_rfl) (boolFin le_rfl) (boolFin le_rfl)

theorem shannonF_not_mem : (fun u => if u ∈ shannonF then (1 : ℝ≥0∞) else 0) ∉ P TT
\end{lstlisting}

\phantomsection\label{lean:tensorgap-bool}
\begin{lstlisting}
abbrev W2 := (W × W) × (W × W)

noncomputable def tens2 (t t' : W × W → ℝ≥0∞) : W2 → ℝ≥0∞ := fun u => t u.1 * t' u.2

def gensTT : Set (W2 → ℝ≥0∞) := {u | ∃ t ∈ tensDual, ∃ t' ∈ tensDual, u = tens2 t t'}

def PTT : Set (W2 → ℝ≥0∞) := orth (orth gensTT)
def PTTdual : Set (W2 → ℝ≥0∞) := orth gensTT

def cohTT : W2 → W2 → Prop := tensC secondC secondC
def cohTTdual : W2 → W2 → Prop := dualC cohTT

def shannon : Finset W2 :=
  {(((true, true), (true, true)), ((true, true), (true, true))),
   (((false, true), (false, true)), ((false, true), (true, false))),
   (((false, true), (true, false)), ((false, false), (false, false))),
   (((true, false), (false, false)), ((false, true), (false, true))),
   (((false, false), (false, false)), ((true, false), (false, false)))}

theorem shannon_clique : IsClique cohTTdual (· ∈ shannon)

theorem shannon_not_mem :
    (fun u => if u ∈ shannon then (1 : ℝ≥0∞) else 0) ∉ PTTdual
\end{lstlisting}

\phantomsection\label{lean:orth-recipe}
\begin{lstlisting}
theorem orth_recipe (A : Fm) : orth (recipe A) = loc (Fm.neg A)

theorem orth_loc (A : Fm) : orth (loc A) = recipe (Fm.neg A)
\end{lstlisting}

\phantomsection\label{lean:tight}
\begin{lstlisting}
theorem tight_iff (A : Fm) :
    (orth (recipe A) = recipe (Fm.neg A) ∧ orth (recipe (Fm.neg A)) = recipe A) ↔
      (recipe A = P A ∧ recipe (Fm.neg A) = P (Fm.neg A))

theorem tight_iff_recipe_eq_loc (A : Fm) :
    orth (recipe A) = recipe (Fm.neg A) ↔ recipe A = loc A

theorem not_tight_core : orth (recipe (core 2 2 2 2)) ≠ recipe (Fm.neg (core 2 2 2 2))
\end{lstlisting}

\phantomsection\label{lean:tight-tensor}
\begin{lstlisting}
theorem P_lolli_eq_loc (n m : ℕ) :
    P (lolli (base n) (base m)) = loc (lolli (base n) (base m))

theorem tight_lolli (n m : ℕ) :
    orth (recipe (lolli (base n) (base m))) = recipe (Fm.neg (lolli (base n) (base m)))

theorem tight_not_tensor_closed {n m n' m' : ℕ} (hn : 2 ≤ n) (hm : 2 ≤ m) (hn' : 2 ≤ n')
    (hm' : 2 ≤ m') :
    orth (recipe (lolli (base n) (base m))) = recipe (Fm.neg (lolli (base n) (base m))) ∧
    orth (recipe (lolli (base n') (base m'))) = recipe (Fm.neg (lolli (base n') (base m'))) ∧
    orth (recipe (tens (lolli (base n) (base m)) (lolli (base n') (base m'))))
      ≠ recipe (Fm.neg (tens (lolli (base n) (base m)) (lolli (base n') (base m'))))
\end{lstlisting}

\phantomsection\label{lean:bell}
\begin{lstlisting}
theorem sharp_dual_exclusive {s : web (core 2 2 2 2) → ℝ≥0∞}
    (hs : s ∈ sharpP (Fm.neg (core 2 2 2 2))) {e e' : BEv} (hne : e ≠ e') (he : s e = 1)
    (he' : s e' = 1) : ¬ cohE e e'

theorem chsh_bilinear {X Y : Fin 2 × Fin 2 → ℝ≥0∞} (hX : ∀ a, ∑ b, X (a, b) ≤ 1)
    (hY : ∀ a, ∑ b, Y (a, b) ≤ 1) :
    ∑ e : BEv, (if win e then X e.1 * Y e.2 else 0) ≤ 3

theorem prBox_mem_loc : prBox ∈ loc (core 2 2 2 2)

theorem prBox_not_mem_P : prBox ∉ P (core 2 2 2 2)

theorem P_core_ssubset_loc : P (core 2 2 2 2) ⊂ loc (core 2 2 2 2)
\end{lstlisting}

\subsection{Quantum Event Logic}

\phantomsection\label{lean:cabello-set}
\begin{lstlisting}
def vec : Fin 18 → Fin 4 → ℤ := ![
  ![-1, 1, 1, 1], ![0, 0, 0, 1], ![0, 0, 1, 0], ![0, 0, 1, 1], ![0, 1, -1, 0], ![0, 1, 0, -1],
  ![0, 1, 0, 0], ![1, -1, -1, 1], ![1, -1, 0, 0], ![1, -1, 1, -1], ![1, 0, -1, 0], ![1, 0, 0, -1],
  ![1, 0, 0, 1], ![1, 0, 1, 0], ![1, 1, -1, 1], ![1, 1, 0, 0], ![1, 1, 1, -1], ![1, 1, 1, 1]]

def bas : Fin 9 → Fin 4 → Fin 18 := ![
  ![1, 2, 15, 8], ![1, 6, 13, 10], ![9, 7, 15, 3], ![9, 17, 10, 5], ![2, 6, 12, 11],
  ![7, 17, 11, 4], ![14, 16, 8, 3], ![14, 0, 13, 5], ![16, 0, 12, 4]]

def dot (u v : Fin 4 → ℤ) : ℤ := ∑ k, u k * v k
\end{lstlisting}

\phantomsection\label{lean:basis-orth}
\begin{lstlisting}
theorem basis_orth : ∀ b : Fin 9, ∀ i j : Fin 4, i ≠ j → dot (vec (bas b i)) (vec (bas b j)) = 0

theorem occ_two : ∀ v : Fin 18,
    (univ.filter (fun p : Fin 9 × Fin 4 => bas p.1 p.2 = v)).card = 2
\end{lstlisting}

\phantomsection\label{lean:no-ks}
\begin{lstlisting}
theorem no_ks : ¬ ∃ f : Fin 18 → Bool, ∀ b, (univ.filter (fun i => f (bas b i))).card = 1
\end{lstlisting}

\phantomsection\label{lean:no-certain-state}
\begin{lstlisting}
def certain : Set (Fin 18 → ℝ≥0∞) :=
  {s | IsSharpVec s ∧ ∀ b, ∑ i, s (bas b i) = 1}

theorem no_certain_state : certain = ∅

theorem Mix_certain_empty : Mix certain = ∅
\end{lstlisting}

\phantomsection\label{lean:parseval}
\begin{lstlisting}
theorem parseval (ψ : Fin 4 → ℝ) (b : Fin 9) :
    ∑ i, (∑ k, ψ k * vec (bas b i) k) ^ 2 / ∑ k, ((vec (bas b i) k : ℝ)) ^ 2 = ∑ k, ψ k ^ 2
\end{lstlisting}

\phantomsection\label{lean:born-not-mix}
\begin{lstlisting}
def states : Set (Fin 18 → ℝ≥0∞) := {x | ∀ b, ∑ i, x (bas b i) = 1}

noncomputable def born (ψ : Fin 4 → ℝ) : Fin 18 → ℝ≥0∞ := fun v =>
  ENNReal.ofReal ((∑ k, ψ k * vec v k) ^ 2 / ∑ k, ((vec v k : ℝ)) ^ 2)

theorem born_mem {ψ : Fin 4 → ℝ} (hψ : ∑ k, ψ k ^ 2 = 1) : born ψ ∈ states

theorem born_not_mix {ψ : Fin 4 → ℝ} (hψ : ∑ k, ψ k ^ 2 = 1) :
    born ψ ∈ states ∧ born ψ ∉ Mix certain
\end{lstlisting}

\phantomsection\label{lean:quantum-partial}
\begin{lstlisting}
def partialCertain : Set (Fin 18 → ℝ≥0∞) :=
  {s | IsSharpVec s ∧ ∀ b, ∑ i, s (bas b i) ≤ 1}

theorem sum_bases (f : Fin 18 → ℝ≥0∞) :
    ∑ b : Fin 9, ∑ i : Fin 4, f (bas b i) = 2 * ∑ v, f v

theorem partial_sum_le_four {s : Fin 18 → ℝ≥0∞} (hs : s ∈ partialCertain) : ∑ v, s v ≤ 4

theorem mix_sum_le_four {x : Fin 18 → ℝ≥0∞} (hx : x ∈ Mix partialCertain) : ∑ v, x v ≤ 4

theorem born_total {ψ : Fin 4 → ℝ} (hψ : ∑ k, ψ k ^ 2 = 1) : 2 * ∑ v, born ψ v = 9

theorem born_not_mix_partial {ψ : Fin 4 → ℝ} (hψ : ∑ k, ψ k ^ 2 = 1) :
    born ψ ∉ Mix partialCertain

def S4 : Finset (Fin 18) := {0, 1, 3, 11}

theorem S4_stable : ∀ b : Fin 9, (Finset.univ.filter (fun i : Fin 4 => bas b i ∈ S4)).card ≤ 1

noncomputable def ind4 : Fin 18 → ℝ≥0∞ := fun v => if v ∈ S4 then 1 else 0

theorem ind4_mem : ind4 ∈ partialCertain ∧ ∑ v, ind4 v = 4
\end{lstlisting}

\phantomsection\label{lean:loc-quantum}
\begin{lstlisting}
def ctx : Set (Finset (Fin 18)) := {t | ∃ b, t = Finset.univ.image (bas b)}

theorem locN_cabello : locN ctx = states

theorem globN_cabello : globN ctx = ∅

theorem born_local {ψ : Fin 4 → ℝ} (hψ : ∑ k, ψ k ^ 2 = 1) : born ψ ∈ locN ctx

theorem born_mono (ψ : Fin 4 → ℝ) : born ψ 1 ≤ born ψ 12 + born ψ 11

def halfSet : Finset (Fin 18) := {0, 1, 2, 6, 7, 9, 14, 16, 17}

noncomputable def half : Fin 18 → ℝ≥0∞ := fun v => if v ∈ halfSet then 2⁻¹ else 0

theorem half_not_born :
    half ∈ locN ctx ∧ half ∉ Mix {x | ∃ ψ : Fin 4 → ℝ, ∑ k, ψ k ^ 2 = 1 ∧ x = born ψ}
\end{lstlisting}

\phantomsection\label{lean:born-ray}
\begin{lstlisting}
noncomputable def unit (u : Fin 18) : Fin 4 → ℝ :=
  fun k => (vec u k : ℝ) / Real.sqrt (dot (vec u) (vec u))

theorem unit_norm (u : Fin 18) : ∑ k, unit u k ^ 2 = 1

theorem born_ray_zero_iff (u v : Fin 18) : born (unit u) v = 0 ↔ dot (vec u) (vec v) = 0

theorem born_ray_one_iff (u v : Fin 18) : born (unit u) v = 1 ↔ u = v
\end{lstlisting}
